% Paper 1 - formal statements (single source). Moved verbatim from sections/04_theory.tex,
% which uses these macros in place.
% T1b (2026-09-28): human-approved hypothesis corrections a-e applied (Corollary 2, Gaussian
% ambiguity cost, Corollary 3 genericity, surrogate-law symbol, onset (A4)); the rendered paper changes.
% Proofs: proofs/<MacroName>.tex.
% Owner: thinker-paper1. Writers use the macros; they do not edit this file.

% --- Theorem 1
\newcommand{\StmtTheoremOne}{\textbf{Theorem 1 (Indistinguishability criterion).} Under the model of Section 3, with common onset time \(\tau\), mechanism instances \((K, \theta_K, b_K)\) and \((F, \theta_F, b_F)\) induce identical distributions over the observation trajectory \(\{y_t\}_{t\ge 0}\) if and only if

\[\theta_K\, C b_K = \theta_F\, C b_F. \tag{$\ast$}\]}

% --- Corollary 1
\newcommand{\StmtCorollaryOne}{\textbf{Corollary 1 (Fault-mimicking compromise --- corrected for degenerate cases).} Fix \(u = Cb_K\), \(v = Cb_F\), and a nonzero compromise severity \(\theta_K \ne 0\). If the fault severity \(\theta_F\) may be any real number, then

\[\exists\, \theta_F \in \mathbb{R} : \theta_K u = \theta_F v \quad\Longleftrightarrow\quad u \in \mathrm{span}(v).\]}

% --- Corollary 2
\newcommand{\StmtCorollaryTwo}{\textbf{Corollary 2 (Separation implies distinguishability, not high accuracy).} If \(\mathrm{rank}\big([\,Cb_K \ \ Cb_F\,]\big) = 2\), no nonzero \(\theta_K,\theta_F\) satisfy (\(\ast\)); for any instance pair with \((\theta_K,\theta_F) \ne (0,0)\), \(K\) and \(F\) are then distinguishable in the population limit.}

% --- Proposition 1 (title)
\newcommand{\StmtPropositionOneTitle}{\textbf{Proposition 1 (Cost of ambiguity is bounded below under confounding --- corrected).}}

% --- Proposition 1 (statement; title and correction note stay in the section)
\newcommand{\StmtPropositionOne}{Suppose (\(\ast\)) holds (the confounded case), \(0 < p < 1\), \(\mathrm{DISRUPT} > 0\), and \(\mathrm{HARM} > C_{OP} \ge 0\) (harm strictly worse than routine containment --- the assumption the earlier version omitted). Then the best achievable expected cost among policies mapping \(\{y_t\}\) to \(\{\text{contain}, \text{leave alone}\}\) is

\[R^\star = \min\{\,p\,C_{OP} + (1-p)\,\mathrm{DISRUPT},\ \ p\,\mathrm{HARM}\,\},\]

the oracle's expected cost is \(R_{\text{oracle}} = p\,C_{OP}\), and

\[\boxed{\mathrm{CoA}^\star \;=\; R^\star - R_{\text{oracle}} \;=\; \min\{\,(1-p)\,\mathrm{DISRUPT},\ \ p\,(\mathrm{HARM} - C_{OP})\,\} \;>\; 0.}\]}

% --- Theorem 2
\newcommand{\StmtTheoremTwo}{\textbf{Theorem 2 (State-feedback confounding requires unobservability).} Under the state-level injection model, let \(\Delta_0 = \theta_K b_K - \theta_F b_F\). Mechanism instances \((K,\theta_K,b_K)\) and \((F,\theta_F,b_F)\) are observationally indistinguishable at every \(t \ge \tau\) if and only if \(\Delta_0\) lies in the unobservable subspace of the pair \((A,C)\), i.e.

\[\Delta_0 \in \bigcap_{k=0}^{n-1} \ker(CA^k),\]

the classical (Kalman) unobservable subspace of \((A,C)\), where \(n\) is the state dimension.}

% --- Corollary 3
\newcommand{\StmtCorollaryThree}{\textbf{Corollary 3 (Theorem 1's confounded pairs are generically state-feedback-fragile).} Theorem 1 only requires \(\Delta_0 \in \ker(C)\) --- a single condition. Theorem 2 requires \(\Delta_0\) to survive \(n\) such conditions simultaneously (\(\ker(C), \ker(CA), \dots, \ker(CA^{n-1})\)). Whenever \((A,C)\) is observable (for fixed \(C \ne 0\), the generic case for a randomly chosen \(A\)), the unobservable subspace is \(\{0\}\), so \textbf{the only state-feedback-robust confounding is the trivial one where \(\theta_K b_K = \theta_F b_F\) exactly as vectors} --- i.e., the two mechanisms must push the identical linear combination of latent state variables, not merely produce the identical \emph{projection} under \(C\).}

% =====================================================================================
% T1 (2026-09-28): definitions and assumptions made explicit.
% Everything below is NEW and additive. The macros above keep rendering exactly the text
% they rendered before; writer-paper1 switches to the macros below in a separate task.
% Convention: \StmtDef... = definition, \StmtAssumption... = assumption,
% \Stmt<Name>Uses = which definitions and assumptions a statement relies on.
% =====================================================================================

% --- Short reference labels (used by the ...Uses macros and by the new statements)
\newcommand{\RefDefOutputModel}{Definition~D1}
\newcommand{\RefDefStateModel}{Definition~D2}
\newcommand{\RefDefInstance}{Definition~D3}
\newcommand{\RefDefIndist}{Definition~D4}
\newcommand{\RefDefPopDist}{Definition~D5}
\newcommand{\RefDefObservability}{Definition~D6}
\newcommand{\RefDefGeneric}{Definition~D7}
\newcommand{\RefDefLossTable}{Definition~D8}
\newcommand{\RefDefDecision}{Definition~D9}
\newcommand{\RefDefTV}{Definition~D10}
\newcommand{\RefDefDiscriminability}{Definition~D11}
\newcommand{\RefAssmNoise}{(A1)}
\newcommand{\RefAssmInit}{(A2)}
\newcommand{\RefAssmCommon}{(A3)}
\newcommand{\RefAssmOnset}{(A4)}
\newcommand{\RefAssmWindow}{(A5)}
\newcommand{\RefAssmPositivity}{(A6)}

% --- Definitions

% D1 --- output-level (sensor-level) injection model (Section 3; used by Theorem 1)
\newcommand{\StmtDefOutputModel}{\textbf{Definition D1 (Output-level injection model).} Fix \(n, m\), a matrix \(A \in \mathbb{R}^{n\times n}\), an observation matrix \(C \in \mathbb{R}^{m\times n}\), covariance matrices \(Q \in \mathbb{R}^{n\times n}\) and \(\Sigma \in \mathbb{R}^{m\times m}\), and an onset time \(\tau\). Under a mechanism instance \((M, \theta_M, b_M)\) (\RefDefInstance), the latent state, the perturbed state and the observation evolve for \(t \ge 0\) as
\[\begin{gathered} x_{t+1} = A x_t + w_t, \qquad \tilde x_t = x_t + \mathbb{1}[t \ge \tau]\,\theta_M\, b_M, \\ y_t = C \tilde x_t + \varepsilon_t, \end{gathered}\]
with \(w_t \sim \mathcal{N}(0,Q)\) and \(\varepsilon_t \sim \mathcal{N}(0,\Sigma)\). The injection enters the observation only; the state recursion does not see it. (Section~3 additionally takes \(m < n\); no statement below uses that restriction.)}

% D2 --- state-level injection model (Section 4 subsection; used by Theorem 2)
\newcommand{\StmtDefStateModel}{\textbf{Definition D2 (State-level injection model).} As \RefDefOutputModel, except that the perturbed state drives the recursion from the onset on:
\[\begin{gathered} x_{t+1} = A x_t + w_t \ (t < \tau), \qquad x_{t+1} = A \tilde x_t + w_t \ (t \ge \tau), \\ \tilde x_t = x_t + \mathbb{1}[t \ge \tau]\,\theta_M\, b_M, \qquad y_t = C \tilde x_t + \varepsilon_t. \end{gathered}\]
The injection is persistent (added at every \(t \ge \tau\)) and propagates through \(A\).}

% D3 --- mechanism instance
\newcommand{\StmtDefMechanismInstance}{\textbf{Definition D3 (Mechanism instance).} A mechanism instance is a triple \((M, \theta_M, b_M)\) with cause \(M \in \{K, F\}\) (\(K\): compromise, \(F\): fault), a deterministic severity \(\theta_M \in \mathbb{R}\) and a deterministic injection direction \(b_M \in \mathbb{R}^n\). Statements compare one fixed instance of \(K\) with one fixed instance of \(F\) (pairwise-instance reading); no distribution over severities, directions or onsets within a cause is modelled.}

% D4 --- observation window, observational indistinguishability
\newcommand{\StmtDefIndistinguishability}{\textbf{Definition D4 (Observation window; observational indistinguishability).} An observation window is a set \(W \subseteq \{0, 1, 2, \dots\}\) of time indices, finite or not; the observation trajectory on \(W\) is \(Y_W = (y_t)_{t \in W}\). Write \(P_M^W\) for the law of \(Y_W\) under instance \((M,\theta_M,b_M)\) (for infinite \(W\), the law on the product \(\sigma\)-algebra, determined by its finite-dimensional marginals). Instances \((K,\theta_K,b_K)\) and \((F,\theta_F,b_F)\) are observationally indistinguishable on \(W\) if \(P_K^W = P_F^W\). Theorem~1 uses \(W = \{t : t \ge 0\}\); ``indistinguishable at every \(t \ge \tau\)'' (Theorem~2) means \(P_K^W = P_F^W\) for \(W = \{t : t \ge \tau\}\); the finite-horizon version of Theorem~2 uses \(W = \{\tau, \dots, \tau+H-1\}\).}

% D5 --- distinguishable in the population limit (Corollary 2)
\newcommand{\StmtDefPopulationDistinguishable}{\textbf{Definition D5 (Distinguishable in the population limit).} Two instances are distinguishable in the population limit on \(W\) if they are not observationally indistinguishable on \(W\), i.e.\ \(P_K^W \ne P_F^W\) (equivalently \(\mathrm{TV}(P_K^W, P_F^W) > 0\), \RefDefTV). This is a statement about the two laws only; it asserts nothing about the accuracy attainable from finitely many observations.}

% D6 --- observability, unobservable subspace (Theorem 2, Corollary 3)
\newcommand{\StmtDefObservability}{\textbf{Definition D6 (Unobservable subspace; observability).} The (Kalman) unobservable subspace of \((A,C)\) is \(\mathcal{N}(A,C) = \bigcap_{k=0}^{n-1} \ker(CA^k)\). The pair \((A,C)\) is observable if \(\mathcal{N}(A,C) = \{0\}\), equivalently if the stacked matrix
\[\begin{bmatrix} C \\ CA \\ \vdots \\ CA^{n-1} \end{bmatrix}\]
has rank \(n\).}

% D7 --- generic (Corollary 3)
\newcommand{\StmtDefGeneric}{\textbf{Definition D7 (Generic).} For fixed \(C\), a property of \(A\) holds generically if it holds for Lebesgue-almost every \(A \in \mathbb{R}^{n\times n}\), i.e.\ it fails only on a Lebesgue-null subset of \(\mathbb{R}^{n\times n}\).}

% D8 --- general loss table and its correspondence with Paper 1's table (human decision 2026-09-28)
\newcommand{\StmtDefLossTable}{\textbf{Definition D8 (Loss table).} The general two-action, two-cause loss table \(L(a, M)\), \(a \in \{\text{contain}, \text{leave alone}\}\), \(M \in \{K,F\}\), is
\[\begin{gathered} L(\text{contain}, K) = C_K, \qquad L(\text{contain}, F) = C_F, \\ L(\text{leave alone}, K) = H_K, \qquad L(\text{leave alone}, F) = H_F, \end{gathered}\]
with \(C_K, C_F, H_K, H_F \in \mathbb{R}\). Paper~1's table (Section~3) is the special case
\[C_{OP} = C_K, \qquad \mathrm{DISRUPT} = C_F, \qquad \mathrm{HARM} = H_K, \qquad H_F = 0.\]
Proposition~1 is stated in Paper~1's symbols \(C_{OP}, \mathrm{DISRUPT}, \mathrm{HARM}\).}

% D9 --- prior, policy class, R*, R_oracle, CoA*
\newcommand{\StmtDefDecisionProblem}{\textbf{Definition D9 (Prior, policies, Bayes risk, oracle risk, cost of ambiguity).} The cause is drawn as \(M = K\) with prior probability \(p = P(K)\) and \(M = F\) with probability \(1-p\); given \(M\), the observation trajectory \(Y_W\) has law \(P_M^W\). A policy is a measurable map \(\pi\) from observation trajectories on \(W\) to \([0,1]\), read as the probability of the action \emph{contain} (randomized policies allowed; a deterministic policy takes values in \(\{0,1\}\)). Its expected cost is
\[\begin{aligned} R(\pi) = {}& p\,\mathbb{E}_K\big[\pi(Y_W)\,C_K + (1-\pi(Y_W))\,H_K\big] \\ &+ (1-p)\,\mathbb{E}_F\big[\pi(Y_W)\,C_F + (1-\pi(Y_W))\,H_F\big]. \end{aligned}\]
The best achievable expected cost is \(R^\star = \inf_\pi R(\pi)\); the oracle, which observes \(M\), has expected cost \(R_{\text{oracle}} = p \min\{C_K, H_K\} + (1-p)\min\{C_F, H_F\}\); the cost of ambiguity is \(\mathrm{CoA}^\star = R^\star - R_{\text{oracle}}\). Where the observation is the raw trajectory \(Y\) the subscript \(Y\) is written (\(R_Y^\star\), \(\mathrm{CoA}_Y^\star\)).}

% D10 --- densities and total variation (TV extension of Proposition 1)
\newcommand{\StmtDefTotalVariation}{\textbf{Definition D10 (Observation densities; total variation).} Let \(\mu\) be a measure dominating \(P_K\) and \(P_F\) (e.g.\ \(P_K + P_F\)), and \(f_K, f_F\) their densities with respect to \(\mu\); \(\int \cdot\, dy\) denotes integration against \(\mu\). The total variation distance is
\[\mathrm{TV}(P,P') = \sup_{B} |P(B) - P'(B)| = \tfrac{1}{2}\int |f - f'|\, dy,\]
so that \(1 - \mathrm{TV}(P_K, P_F) = \int \min\{f_K, f_F\}\, dy\).}

% D11 --- Gaussian discriminability (Corollary 2 benchmark, CoA_Y*(d))
\newcommand{\StmtDefDiscriminability}{\textbf{Definition D11 (Gaussian discriminability).} If \(Y \mid M \sim \mathcal{N}(\mu_M, V)\), \(M \in \{K,F\}\), with a common covariance \(V \succ 0\), let \(\delta = \mu_K - \mu_F\) and \(d = \sqrt{\delta^\top V^{-1} \delta} \ge 0\) (the whitened mean separation).}

% --- Assumptions (numbered list; each statement's ...Uses macro references the items it needs)
\newcommand{\StmtAssmNoise}{\textbf{\RefAssmNoise{} Gaussian, independent noise.} \(w_t \sim \mathcal{N}(0,Q)\) and \(\varepsilon_t \sim \mathcal{N}(0,\Sigma)\), with \(Q\) and \(\Sigma\) symmetric positive semidefinite; the family \(\{w_t\}_{t \ge 0} \cup \{\varepsilon_t\}_{t \ge 0}\) is mutually independent (across time and across the two sequences), and its law does not depend on the instance.}
\newcommand{\StmtAssmInit}{\textbf{\RefAssmInit{} Deterministic initial state.} \(x_0 = 0\). (The proofs use only that \(x_0\) is Gaussian, possibly degenerate, independent of the noise, with a law that does not depend on the instance; \(x_0 = 0\) is the implemented special case.)}
\newcommand{\StmtAssmCommon}{\textbf{\RefAssmCommon{} Known, common model.} \(A\), \(Q\), \(C\), \(\Sigma\) are known and identical under both causes; only the injection term \(\theta_M b_M\) depends on the instance.}
\newcommand{\StmtAssmOnset}{\textbf{\RefAssmOnset{} Common onset.} The onset \(\tau\) is deterministic, common to both causes, not necessarily known to the operator. (The proofs use only that \(\tau\) is deterministic and common; no statement requires a policy to know \(\tau\).)}
\newcommand{\StmtAssmWindow}{\textbf{\RefAssmWindow{} Window reaches the onset.} The observation window \(W\) contains at least one \(t \ge \tau\). Theorem~2 uses the stronger forms \(W \supseteq \{t : t \ge \tau\}\) (all-time version) or \(W \supseteq \{\tau, \dots, \tau+H-1\}\) (finite-horizon version).}
\newcommand{\StmtAssmPositivity}{\textbf{\RefAssmPositivity{} Positive-definite observation noise.} \(\Sigma \succ 0\). Used only to apply the Gaussian ambiguity cost to the model on a finite window: it makes the covariance of the stacked observations positive definite (\(V \succ 0\)). No other statement uses a positivity condition on \(\Sigma\) or \(Q\); \(Q\) is only required to be positive semidefinite.}
\newcommand{\StmtAssumptionsList}{\textbf{Standing assumptions.}
\begin{enumerate}
\item[] \StmtAssmNoise
\item[] \StmtAssmInit
\item[] \StmtAssmCommon
\item[] \StmtAssmOnset
\item[] \StmtAssmWindow
\item[] \StmtAssmPositivity
\end{enumerate}}

% --- What each existing statement relies on (reference macros; the statements above are unchanged)
\newcommand{\StmtTheoremOneUses}{\emph{Theorem~1 uses} \RefDefOutputModel, \RefDefInstance, \RefDefIndist{} (window \(W = \{t: t\ge 0\}\)) and assumptions \RefAssmNoise--\RefAssmWindow.}
\newcommand{\StmtCorollaryOneUses}{\emph{Corollary~1 uses} \RefDefInstance{} and condition (\(\ast\)) of Theorem~1; the equivalence itself is linear algebra and uses no probabilistic assumption.}
\newcommand{\StmtCorollaryTwoUses}{\emph{Corollary~2 uses} \RefDefInstance, \RefDefIndist, \RefDefPopDist{} and, through Theorem~1, \RefDefOutputModel{} and assumptions \RefAssmNoise--\RefAssmWindow.}
\newcommand{\StmtPropositionOneUses}{\emph{Proposition~1 uses} \RefDefLossTable{} (Paper~1's symbols), \RefDefDecision{} and, through Theorem~1 (to pass from (\(\ast\)) to \(P_K^W = P_F^W\)), \RefDefOutputModel, \RefDefInstance, \RefDefIndist{} and assumptions \RefAssmNoise--\RefAssmWindow.}
\newcommand{\StmtPropositionOneTVUses}{\emph{The total-variation extension uses} \RefDefLossTable{} (Paper~1's symbols), \RefDefDecision{} and \RefDefTV; it uses none of \RefAssmNoise--\RefAssmPositivity: it holds for any two observation laws.}
\newcommand{\StmtGaussianAmbiguityCostUses}{\emph{The Gaussian ambiguity cost uses} \RefDefLossTable{} (general symbols), \RefDefDecision, \RefDefTV{} and \RefDefDiscriminability; applied to the model on a finite window it needs \RefAssmNoise--\RefAssmWindow{} (common covariance) and \RefAssmPositivity{} (\(V \succ 0\)).}
\newcommand{\StmtTheoremTwoUses}{\emph{Theorem~2 uses} \RefDefStateModel, \RefDefInstance, \RefDefIndist{} (window \(W = \{t : t \ge \tau\}\)), \RefDefObservability{} and assumptions \RefAssmNoise--\RefAssmOnset{} with \RefAssmWindow{} in its all-time form.}
\newcommand{\StmtTheoremTwoFiniteUses}{\emph{The finite-horizon version of Theorem~2 uses} \RefDefStateModel, \RefDefInstance, \RefDefIndist{} (window \(W = \{\tau,\dots,\tau+H-1\}\)), \RefDefObservability{} and assumptions \RefAssmNoise--\RefAssmOnset{} with \RefAssmWindow{} in its finite-horizon form.}
\newcommand{\StmtCorollaryThreeUses}{\emph{Corollary~3 uses} everything Theorem~2 uses, and \RefDefObservability; its accompanying remark also uses \RefDefGeneric.}

% --- Promoted result (a): total-variation extension of Proposition 1 (content moved from the prose of sections/04_theory.tex, unchanged)
\newcommand{\StmtPropositionOneTV}{\textbf{Proposition~1 (total-variation extension).} Let the loss table and prior be as in Proposition~1, with \(0 < p < 1\), \(\mathrm{DISRUPT} > 0\) and \(\mathrm{HARM} > C_{OP}\), but do not assume (\(\ast\)). Let \(f_K, f_F\) be the observation densities (\RefDefTV) and set \(a = p(\mathrm{HARM}-C_{OP})\), \(b = (1-p)\,\mathrm{DISRUPT}\). Then
\[\mathrm{CoA}^\star = \int \min\{a f_K(y),\, b f_F(y)\}\, dy = \frac{a+b-\int |af_K - bf_F|\,dy}{2},\]
\[\min(a,b)\,\big(1-\mathrm{TV}(P_K,P_F)\big) \;\le\; \mathrm{CoA}^\star \;\le\; \min(a,b),\]
and, if a common surrogate law \(\bar P\) satisfies \(\mathrm{TV}(P_K,\bar P) \le \epsilon_K\) and \(\mathrm{TV}(P_F,\bar P) \le \epsilon_F\), then
\[\mathrm{CoA}^\star \ge \min(a,b)\max\{0,\ 1-\epsilon_K-\epsilon_F\}.\]}

% --- Promoted result (b): finite-horizon version of Theorem 2 (content moved from the prose of sections/04_theory.tex, unchanged)
\newcommand{\StmtTheoremTwoFinite}{\textbf{Theorem~2 (finite-horizon version).} Under the state-level injection model, let \(\Delta_0 = \theta_K b_K - \theta_F b_F\), let \(H \ge 1\) be the number of post-onset observations, and define
\[G_H = \begin{bmatrix} C \\ C(I+A) \\ \vdots \\ C\sum_{j=0}^{H-1}A^j \end{bmatrix}.\]
Mechanism instances \((K,\theta_K,b_K)\) and \((F,\theta_F,b_F)\) induce the same law of the observations over the window of \(H\) post-onset observations if and only if \(G_H\Delta_0 = 0\), which holds if and only if \(CA^j\Delta_0 = 0\) for \(j = 0, \dots, H-1\). If \(H \ge n\), this is equivalent to the all-time condition of Theorem~2, \(\Delta_0 \in \bigcap_{k=0}^{n-1}\ker(CA^k)\).}

% --- Promoted result (c): Gaussian unequal-prior ambiguity cost (content moved from the prose of sections/04_theory.tex, unchanged)
\newcommand{\StmtGaussianAmbiguityCost}{\textbf{Proposition (Gaussian unequal-prior ambiguity cost).} Let \(Y \mid M \sim \mathcal{N}(\mu_M, V)\), \(M \in \{K,F\}\), with common covariance \(V \succ 0\) and discriminability \(d > 0\) (\RefDefDiscriminability), prior \(p = P(K)\), and the general loss table with \(a = p(H_K-C_K)\), \(b=(1-p)(C_F-H_F)\), \(a, b > 0\). Then the population-optimal ambiguity cost is
\[\mathrm{CoA}_Y^\star(d) = a\,\Phi\!\left(\frac{\log(b/a)}{d} - \frac{d}{2}\right) + b\,\Phi\!\left(\frac{\log(a/b)}{d} - \frac{d}{2}\right),\]
which tends to \(\min(a,b)\) as \(d \to 0^+\) (the value at \(d=0\), exact confounding, recovering Proposition~1, is obtained as this limit) and decreases toward \(0\) as \(d\to\infty\).}

% --- Corollary 3 split (human decision 2026-09-28). \StmtCorollaryThree above is the legacy combined wording (correction c applied to both).
\newcommand{\StmtCorollaryThreePrecise}{\textbf{Corollary 3 (State-feedback confounding under observability).} Under the hypotheses of Theorem~2, if \((A,C)\) is observable, then mechanism instances \((K,\theta_K,b_K)\) and \((F,\theta_F,b_F)\) are observationally indistinguishable at every \(t \ge \tau\) if and only if \(\theta_K b_K = \theta_F b_F\).}
\newcommand{\StmtCorollaryThreeRemark}{\textbf{Remark (Theorem 1's confounded pairs are generically state-feedback-fragile).} Theorem 1 only requires \(\Delta_0 \in \ker(C)\) --- a single condition. Theorem 2 requires \(\Delta_0\) to survive \(n\) such conditions simultaneously (\(\ker(C), \ker(CA), \dots, \ker(CA^{n-1})\)). Whenever \((A,C)\) is observable (for fixed \(C \ne 0\), the generic case for a randomly chosen \(A\)), the unobservable subspace is \(\{0\}\), so \textbf{the only state-feedback-robust confounding is the trivial one where \(\theta_K b_K = \theta_F b_F\) exactly as vectors} --- i.e., the two mechanisms must push the identical linear combination of latent state variables, not merely produce the identical \emph{projection} under \(C\).}
